\thispagestyle{empty}
\begin{center}
{\huge\bfseries Abstract}
\end{center}
\initial{I}oT-REAP (Internet of Things - Remote Equipment Access Platform)
is a secure web platform for industrial remote laboratory operations, developed
as an end-of-study internship at Novation City, a Tunisian Industry~4.0
cluster. It addresses gaps in solutions - client installation dependency,
absent VM lifecycle management, fragmented tooling, and no integrated training
- within a browser-native interface.

Built on Laravel~12, React~18, TypeScript, and Inertia.js, it integrates
Proxmox~VE for VM provisioning, Apache Guacamole for clientless RDP/VNC/SSH
access, and camera integration with
camera feeds via WebRTC/HLS through MediaMTX. An e-learning module covers
training paths, video lessons, quizzes, and certificates. Security relies on
zero-trust, comprehensive audit logging, Purdue IT/OT isolation, and HMAC webhook
validation; OWASP~ZAP found no critical vulnerabilities. Deployment runs via
GitHub Actions CI/CD to OVH.

\vspace{0.5cm}
\noindent\textbf{Keywords:} Industry~4.0, Remote Laboratory, Virtual Machine Provisioning, Apache Guacamole, MQTT, Zero-Trust Security, E-Learning

\vspace{0.9cm}
\begin{center}
{\huge\bfseries Résumé}
\end{center}
\initial{I}oT-REAP (Internet of Things - Plateforme d'Accès aux Équipements
à Distance) est une plateforme web sécurisée pour les opérations industrielles
à distance, développée lors d'un stage de fin de études à Novation City,
cluster tunisien orienté Industrie~4.0. Elle répond aux lacunes des solutions
existantes - dépendance à l'installation client, absence de gestion du cycle
de vie des machines virtuelles, outils fragmentés et absence de formation
intégrée - au sein d'une interface native navigateur.

Développée avec Laravel~12, React~18, TypeScript et Inertia.js, la plateforme
intègre Proxmox~VE pour le provisionnement de VM, Apache Guacamole pour
l'accès bureau à distance sans client via RDP/VNC/SSH, et Mosquitto~MQTT pour
le contrôle des caméras. Les flux caméra transitent via
WebRTC/HLS par MediaMTX. Un module e-learning propose des parcours de
formation, des leçons vidéo, des quiz et des certificats vérifiables. La
sécurité repose sur zero-trust, un audit complet, l'isolation IT/OT selon
Purdue, et la validation HMAC des webhooks~; OWASP~ZAP n'a révélé aucune
vulnérabilité critique. Le déploiement est automatisé via GitHub Actions
CI/CD vers un serveur OVH.

\vspace{0.5cm}
\noindent\textbf{Mots-clés~:} Industrie~4.0, Laboratoire Distant, Provisionnement de Machines Virtuelles, Apache Guacamole, MQTT, Sécurité Zero-Trust, E-Learning
\cleardoublepage\cleardoublepage